\section{原子核・放射線}

\subsection{平成23年度 問10}
\label{q:h23-reg-applied-10}

\noindent\textbf{原典：} sources/H23/17\_applied\_science.pdf、PDF物理4頁。

\textbf{公式正答：} 2（single）。

\par\textbf{検討未確定：} 原典3は同質量数の最小質量核へβ壊変するという全称。単一β安定核が複数となる偶数Aの対エネルギーを省く。二重βの実在は単一β経路とは別である。原典4は内部取込みの量・経路条件を省く。公式・初回2を保持し、3の科学的排他性は未確定。

\paragraph{問題の要旨} 天然壊変・熱中性子核分裂・β安定・透過と内部影響・γ脱励起を区別する。

\Needspace{4\baselineskip}\paragraph{確認した条件・記述（独自要約）}\mbox{}\par

\noindent\textbf{方向：} 誤っているもの。\par

\noindent\textbf{対象：} 量子・原子核の性質。質量数、係数、否定と肢対応を画像照合。\par

\paragraph{問われる核心} 核反応のエネルギー条件と保存するA・Z

\paragraph{必要知識} αでA-4/Z-2、βでA一定/Z±1、γで核種不変、熱核分裂と高速核分裂、質量放物線と対エネルギー、内部と外部の透過条件。

\paragraph{考え方・正答を選ぶ理由（査読待ち初稿）} 2の238Uと熱中性子が取り違えである。235U等の熱核分裂、238Uの高速核分裂・中性子捕獲を分ける。3はβ安定へ質量が減る一般傾向を述べるが、偶奇性と隣接核のQ値を省いた全称では成り立たない。公式2・独立2を保持し3の文字通りの一意性は査読対象に残す。

\Needspace{4\baselineskip}\paragraph{選択肢の検討（査読待ち初稿）}\mbox{}\par

\noindent\textbf{1：主要天然系列で適切} 自然UとThはα・β壊変を経て安定Pbへ至る。。主要な天然U238/U235とTh232の壊変系列の終点は安定Pb同位体。原子番号を変えるα・βと核準位遷移を分ける。\par

\noindent\textbf{2：不適切} 238Uへ熱中性子を当て核分裂連鎖反応で発電。。通常の熱中性子連鎖反応では235U等の核分裂性核種を用いる。238Uは高速中性子核分裂や捕獲後Pu生成の役割を持ち、熱核分裂性とは区別する。\par

\noindent\textbf{3：一般傾向・全称疑義} 同質量数の核はその中の最小質量へ向けてβ壊変。。β壊変はエネルギー的に可能な隣接Zへ進む。偶数Aでは偶偶・奇奇質量が別放物線となり、単一βに安定な核が複数あり得る。最小質量へ必ず単一βで進む全称は不正確。\par

\noindent\textbf{4：適切（影響条件付き）} α線は紙で止まるが、放出核種を体内へ取り込むとダメージが大きい。。外部αの透過は弱いが内部取込みでは短い距離に高密度にエネルギーを与え得る。影響は量・核種・摂取経路・組織への分布等で変わるので無条件に一定被害を割り当てない。\par

\noindent\textbf{5：適切} γ放出はエネルギー準位が変わり原子番号は変わらない。。γは同じ核種の励起準位からの電磁遷移で、AとZを保つ。α・β後に伴うγは先行壊変で変わったZと区別する。\par

\paragraph{再利用する考え方} 定義の基準日を確定し、物理量の関係から指数と係数を導く。量子現象では運動量・エネルギー保存と観測対象を区別する。

\paragraph{共通理論} \hyperref[theory:modern:nuclear]{第\ref*{theory:modern:nuclear}節}, \hyperref[theory:modern:stability]{第\ref*{theory:modern:stability}節}, \hyperref[theory:modern:fission]{第\ref*{theory:modern:fission}節}, \hyperref[theory:modern:dose]{第\ref*{theory:modern:dose}節}

\par 参考：\cite{modern_mit_nuclear,modern_mit_mass,modern_doublebeta,modern_fastfission,modern_alpha,modern_series}。

\par 難易度 2（校準初稿）、優先度未評価。

\subsection{平成26年度 問9}
\label{q:h26-reg-applied-09}

\noindent\textbf{原典：} sources/H26/17\_applied\_science.pdf、PDF物理4頁。

\textbf{公式正答：} 2（single）。

\paragraph{問題の要旨} 半減期と経過時間の比から残存する親核種の割合を求める。

\Needspace{4\baselineskip}\paragraph{確認した条件・記述（独自要約）}\mbox{}\par

\noindent\textbf{方向：} 最も適切。\par

\noindent\textbf{半減期：} 30年\par

\noindent\textbf{経過時間：} 90年\par

\noindent\textbf{対象：} 元の放射性核種の残存量、生成等なし。\par

\paragraph{問われる核心} 半減期は指数減衰の時間尺度

\paragraph{必要知識} 一定壊変定数、T=\allowbreak{}ln2/λ、N/N0=\allowbreak{}2\textasciicircum{}(-t/T)、親核種量、指定時間の同単位化。

\paragraph{考え方・正答を選ぶ理由（査読待ち初稿）} 生成・流出入を除き一定壊変率なら\(dN/dt=-\lambda N\)を積分して\(N/N_0=e^{-\lambda t}=2^{-t/T_{1/2}}\)。\(t/T_{1/2}=90/30\)から残存比は0.125、2の1/8を選ぶ。残った親核種量と壊変済み割合1−N/N0を区別する。

\Needspace{4\baselineskip}\paragraph{選択肢の検討（査読待ち初稿）}\mbox{}\par

\noindent\textbf{1：不一致} 1/6。N/N0=\allowbreak{}2\textasciicircum{}(-t/T)で得る0.125と比較すると1/6は大きすぎる。 時間の比を線形に割るのではない。\par

\noindent\textbf{2：適切} 1/8。計算結果と一致する。\par

\noindent\textbf{3：不一致} 1/9。N/N0=\allowbreak{}2\textasciicircum{}(-t/T)で得る0.125と比較すると1/9は小さすぎる。 時間の比を線形に割るのではない。\par

\noindent\textbf{4：不一致} 1/16。N/N0=\allowbreak{}2\textasciicircum{}(-t/T)で得る0.125と比較すると1/16は小さすぎる。 時間の比を線形に割るのではない。\par

\noindent\textbf{5：不一致} 1/18。N/N0=\allowbreak{}2\textasciicircum{}(-t/T)で得る0.125と比較すると1/18は小さすぎる。 時間の比を線形に割るのではない。\par

\paragraph{再利用する考え方} 定義の基準日を確定し、物理量の関係から指数と係数を導く。量子現象では運動量・エネルギー保存と観測対象を区別する。

\paragraph{共通理論} \hyperref[theory:modern:decay]{第\ref*{theory:modern:decay}節}

\par 参考：\cite{modern_mit_nuclear}。

\paragraph{関連問題} \hyperref[q:h28-reg-applied-09]{H28-REG-APPLIED-09}

\par 難易度 1（校準初稿）、優先度未評価。

\subsection{平成28年度 問9}
\label{q:h28-reg-applied-09}

\noindent\textbf{原典：} sources/H28/17\_applied\_science.pdf、PDF物理4頁。

\textbf{公式正答：} 3（single）。

\paragraph{問題の要旨} 半減期と経過時間の比から残存する親核種の割合を求める。

\Needspace{4\baselineskip}\paragraph{確認した条件・記述（独自要約）}\mbox{}\par

\noindent\textbf{方向：} 最も近い値。\par

\noindent\textbf{半減期：} 8日\par

\noindent\textbf{経過時間：} 30日\par

\noindent\textbf{対象：} 元の放射性核種の残存量、生成等なし。\par

\paragraph{問われる核心} 半減期は指数減衰の時間尺度

\paragraph{必要知識} 一定壊変定数、T=\allowbreak{}ln2/λ、N/N0=\allowbreak{}2\textasciicircum{}(-t/T)、親核種量、指定時間の同単位化。

\paragraph{考え方・正答を選ぶ理由（査読待ち初稿）} 生成・流出入を除き一定壊変率なら\(dN/dt=-\lambda N\)を積分して\(N/N_0=e^{-\lambda t}=2^{-t/T_{1/2}}\)。\(t/T_{1/2}=30/8\)から残存比は0.07432544、3の.07を選ぶ。残った親核種量と壊変済み割合1−N/N0を区別する。

\Needspace{4\baselineskip}\paragraph{選択肢の検討（査読待ち初稿）}\mbox{}\par

\noindent\textbf{1：不一致} .20。N/N0=\allowbreak{}2\textasciicircum{}(-t/T)で得る0.0743254と比較すると.20は大きすぎる。 適切肢より差が大きい。\par

\noindent\textbf{2：不一致} .15。N/N0=\allowbreak{}2\textasciicircum{}(-t/T)で得る0.0743254と比較すると.15は大きすぎる。 適切肢より差が大きい。\par

\noindent\textbf{3：適切} .07。計算値へ5候補中で最も近い。\par

\noindent\textbf{4：不一致} .05。N/N0=\allowbreak{}2\textasciicircum{}(-t/T)で得る0.0743254と比較すると.05は小さすぎる。 適切肢より差が大きい。\par

\noindent\textbf{5：不一致} .03。N/N0=\allowbreak{}2\textasciicircum{}(-t/T)で得る0.0743254と比較すると.03は小さすぎる。 適切肢より差が大きい。\par

\paragraph{再利用する考え方} 定義の基準日を確定し、物理量の関係から指数と係数を導く。量子現象では運動量・エネルギー保存と観測対象を区別する。

\paragraph{共通理論} \hyperref[theory:modern:decay]{第\ref*{theory:modern:decay}節}

\par 参考：\cite{modern_mit_nuclear}。

\paragraph{関連問題} \hyperref[q:h26-reg-applied-09]{H26-REG-APPLIED-09}, \hyperref[q:r04-reg-applied-09]{R04-REG-APPLIED-09}

\par 難易度 1（校準初稿）、優先度未評価。

\subsection{令和4年度 問9}
\label{q:r04-reg-applied-09}

\noindent\textbf{原典：} sources/R04/17\_applied\_science.pdf、PDF物理3頁。

\textbf{公式正答：} 2（single）。

\paragraph{問題の要旨} 半減期と経過時間の比から残存する親核種の割合を求める。

\Needspace{4\baselineskip}\paragraph{確認した条件・記述（独自要約）}\mbox{}\par

\noindent\textbf{方向：} 最も近い値。\par

\noindent\textbf{半減期：} 9日\par

\noindent\textbf{経過時間：} 30日\par

\noindent\textbf{対象：} 元の放射性核種の残存量、生成等なし。\par

\paragraph{問われる核心} 半減期は指数減衰の時間尺度

\paragraph{必要知識} 一定壊変定数、T=\allowbreak{}ln2/λ、N/N0=\allowbreak{}2\textasciicircum{}(-t/T)、親核種量、指定時間の同単位化。

\paragraph{考え方・正答を選ぶ理由（査読待ち初稿）} 生成・流出入を除き一定壊変率なら\(dN/dt=-\lambda N\)を積分して\(N/N_0=e^{-\lambda t}=2^{-t/T_{1/2}}\)。\(t/T_{1/2}=30/9\)から残存比は0.09921257、2の.10を選ぶ。残った親核種量と壊変済み割合1−N/N0を区別する。

\Needspace{4\baselineskip}\paragraph{選択肢の検討（査読待ち初稿）}\mbox{}\par

\noindent\textbf{1：不一致} .20。N/N0=\allowbreak{}2\textasciicircum{}(-t/T)で得る0.0992126と比較すると.20は大きすぎる。 適切肢より差が大きい。\par

\noindent\textbf{2：適切} .10。計算値へ5候補中で最も近い。\par

\noindent\textbf{3：不一致} .05。N/N0=\allowbreak{}2\textasciicircum{}(-t/T)で得る0.0992126と比較すると.05は小さすぎる。 適切肢より差が大きい。\par

\noindent\textbf{4：不一致} .03。N/N0=\allowbreak{}2\textasciicircum{}(-t/T)で得る0.0992126と比較すると.03は小さすぎる。 適切肢より差が大きい。\par

\noindent\textbf{5：不一致} .01。N/N0=\allowbreak{}2\textasciicircum{}(-t/T)で得る0.0992126と比較すると.01は小さすぎる。 適切肢より差が大きい。\par

\paragraph{再利用する考え方} 定義の基準日を確定し、物理量の関係から指数と係数を導く。量子現象では運動量・エネルギー保存と観測対象を区別する。

\paragraph{共通理論} \hyperref[theory:modern:decay]{第\ref*{theory:modern:decay}節}

\par 参考：\cite{modern_mit_nuclear}。

\paragraph{関連問題} \hyperref[q:h28-reg-applied-09]{H28-REG-APPLIED-09}

\par 難易度 1（校準初稿）、優先度未評価。

\section{光と電子励起を利用する素子}

\subsection{令和5年度 問7}
\label{q:r05-reg-applied-07}

\noindent\textbf{原典：} sources/R05/17\_applied\_science.pdf、PDF物理3頁。

\textbf{公式正答：} 4（single）。

\paragraph{問題の要旨} 光と電子の励起・再結合を利用する素子と、超伝導の巨視的量子現象を比較する。

\Needspace{4\baselineskip}\paragraph{確認した条件・記述（独自要約）}\mbox{}\par

\noindent\textbf{方向：} 最も不適切なもの。\par

\noindent\textbf{対象：} 量子・原子核の性質。質量数、係数、否定と肢対応を画像照合。\par

\paragraph{問われる核心} 光素子の吸収・自発放出・誘導放出の役割

\paragraph{必要知識} 光子E=\allowbreak{}hf、半導体の電子正孔、分子励起、光吸収と放射再結合、レーザーの反転分布/共振器、超伝導の凝縮状態。

\paragraph{考え方・正答を選ぶ理由（査読待ち初稿）} 4がこの分類で最も不適切である。1は光吸収で電子励起を作り電流へ、2・3・5は励起された電子状態の放射緩和・誘導放出で光へ変換する。一方、超伝導は低温の凝縮量子状態の電気磁気応答が主原理である。光励起で超伝導を変化させる研究が存在し得ることとは別の比較である。

\Needspace{4\baselineskip}\paragraph{選択肢の検討（査読待ち初稿）}\mbox{}\par

\noindent\textbf{1：該当} 太陽電池。光吸収で電子・正孔を作り電流として取り出す。光から電気への変換であり光電子相互作用を利用する。\par

\noindent\textbf{2：該当} 発光ダイオードLED。注入した電子と正孔の放射再結合で光を出す。電気から光への変換である。\par

\noindent\textbf{3：該当} 半導体レーザー。活性層の電子励起と反転分布、光電磁場による誘導放出・共振器帰還を利用する。\par

\noindent\textbf{4：主原理が異なる} 超伝導（超電導）。超伝導は電子対等の巨視的な量子状態に伴う無抵抗・磁場応答の現象である。光吸収/放出素子とは主原理が異なる。光と超伝導体が相互作用しないという断言ではない。\par

\noindent\textbf{5：該当} 有機EL。有機半導体へ電荷を注入し、励起子等の放射緩和で光を出す。分子励起も電子励起に含む。\par

\paragraph{再利用する考え方} 定義の基準日を確定し、物理量の関係から指数と係数を導く。量子現象では運動量・エネルギー保存と観測対象を区別する。

\paragraph{共通理論} \hyperref[theory:modern:devices]{第\ref*{theory:modern:devices}節}, \hyperref[theory:modern:superconduct]{第\ref*{theory:modern:superconduct}節}, \hyperref[theory:modern:duality]{第\ref*{theory:modern:duality}節}

\par 参考：\cite{modern_pv,modern_led,modern_laser,modern_oled,modern_superconduct,modern_pairs}。

\par 難易度 2（校準初稿）、優先度未評価。

\section{光子の量子性・電子の物質波}

\subsection{平成24年度 問7}
\label{q:h24-reg-applied-07}

\noindent\textbf{原典：} sources/H24/17\_applied\_science.pdf、PDF物理2頁。

\textbf{公式正答：} 2（single）。

\paragraph{問題の要旨} 光の粒子性に直接関係しない現象又は概念を選ぶ。

\Needspace{4\baselineskip}\paragraph{確認した条件・記述（独自要約）}\mbox{}\par

\noindent\textbf{方向：} 光の粒子性に直接関係しないことがらとして最適。\par

\paragraph{問われる核心} 光子の粒子性と電子の物質波

\paragraph{必要知識} 光子の粒子性と電子の物質波。光子のエネルギーと運動量に関する1・3・4・5に対し、2は電子の波動性を観測する現象である。光の粒子性と物質粒子の波動性を区別して2を選ぶ。

\paragraph{考え方・正答を選ぶ理由（査読待ち初稿）} 光子のエネルギーと運動量に関する1・3・4・5に対し、2は電子の波動性を観測する現象である。光の粒子性と物質粒子の波動性を区別して2を選ぶ。

\Needspace{4\baselineskip}\paragraph{選択肢の検討（査読待ち初稿）}\mbox{}\par

\noindent\textbf{1：直接関係} 光電効果。光がエネルギーhf単位で電子へ渡すモデルで、閾値周波数と光電子の最大運動エネルギーを説明する。\par

\noindent\textbf{2：直接関係しない} 電子線回折。電子の物質波が結晶で干渉・回折する現象で、光の粒子性を直接示すものではない。\par

\noindent\textbf{3：直接関係} コンプトン効果。光子のエネルギー・運動量を電子との散乱で保存させ、波長変化を説明する。\par

\noindent\textbf{4：光電効果へ関係} 仕事関数。表面から電子を取り出す最小エネルギーで、光電効果のエネルギー収支へ入る。これ自体を光子と同一の概念にはしない。\par

\noindent\textbf{5：直接関係} プランク定数。光子エネルギーhfや運動量h/λを結ぶ定数で、光の量子性へ直接関係する。\par

\paragraph{再利用する考え方} 光子の粒子性と電子の物質波を暗記だけでなく、位相・幾何・符号と成立条件から導く。選択肢番号を類似年度へそのまま移さない。

\paragraph{共通理論} \hyperref[theory:modern:duality]{第\ref*{theory:modern:duality}節}

\par 参考：\cite{os_photoelectric,os_matter}。

\par 難易度 2（校準初稿）、優先度未評価。

\subsection{平成25年度 問7}
\label{q:h25-reg-applied-07}

\noindent\textbf{原典：} sources/H25/17\_applied\_science.pdf、PDF物理4頁。

\textbf{公式正答：} 5（single）。

\paragraph{問題の要旨} 同じ速さを持つ粒子の物質波長を質量比から並べる。

\Needspace{4\baselineskip}\paragraph{確認した条件・記述（独自要約）}\mbox{}\par

\noindent\textbf{方向：} 長い順番として適切。\par

\noindent\textbf{対象：} 量子・原子核の性質。質量数、係数、否定と肢対応を画像照合。\par

\paragraph{問われる核心} 物質波長を運動量で比較する

\paragraph{必要知識} λ=h/p、同速のγ共通、電子・陽子・α核の静止質量の順、速度ゼロで波長比較をする極限は除く。

\paragraph{考え方・正答を選ぶ理由（査読待ち初稿）} 5である。\(\lambda=h/p,\quad p=\gamma mv\)。同じ速さではγとvは共通なので\(\lambda\propto1/m\)、電子、陽子、αの順となる。非相対論近似ならp=\allowbreak{}mvで同じ結論であり、同運動エネルギーという別条件を混ぜない。

\Needspace{4\baselineskip}\paragraph{選択肢の検討（査読待ち初稿）}\mbox{}\par

\noindent\textbf{1：誤順序} 陽子→α→電子。電子は最軽量で最長波長なので最後にはならない。\par

\noindent\textbf{2：誤順序} α→陽子→電子。質量の大きい順を波長の大きい順へ取り違えている。\par

\noindent\textbf{3：誤順序} 電子→α→陽子。電子は先頭だがαは陽子より重く、αの波長は短い。\par

\noindent\textbf{4：誤順序} 陽子→電子→α。陽子より軽い電子が最長であり最初の2者が逆。\par

\noindent\textbf{5：適切} 電子→陽子→α。静止質量は電子<陽子<α。同速ではγも共通、p=γmvは質量に比例、λはその逆なのでこの順になる。\par

\paragraph{再利用する考え方} 定義の基準日を確定し、物理量の関係から指数と係数を導く。量子現象では運動量・エネルギー保存と観測対象を区別する。

\paragraph{共通理論} \hyperref[theory:modern:matter]{第\ref*{theory:modern:matter}節}

\par 参考：\cite{os_matter}。

\par 難易度 2（校準初稿）、優先度未評価。

\subsection{平成27年度 問5}
\label{q:h27-reg-applied-05}

\noindent\textbf{原典：} sources/H27/17\_applied\_science.pdf、PDF物理2頁。

\textbf{公式正答：} 2（single）。

\paragraph{問題の要旨} 加速電子の運動エネルギーから連続X線の最短波長を求める。

\Needspace{4\baselineskip}\paragraph{確認した条件・記述（独自要約）}\mbox{}\par

\noindent\textbf{方向：} 最短波長に最も近い値。\par

\noindent\textbf{加速電圧：} 35kV\par

\noindent\textbf{h：} \(6.6\times10^{-34}\mathrm{J\,s}\)\par

\noindent\textbf{e：} \(1.6\times10^{-19}\mathrm C\)\par

\noindent\textbf{c：} \(3.0\times10^8\mathrm{m/s}\)\par

\paragraph{問われる核心} 制動放射の端をエネルギー保存で求める

\paragraph{必要知識} 電圧仕事eV、光子E=\allowbreak{}hc/λ、全エネルギーを一光子へ渡す最短端、kV=\allowbreak{}10\textsuperscript{3}V、指定h/e/c。

\paragraph{考え方・正答を選ぶ理由（査読待ち初稿）} 初期運動エネルギーを無視する電子が電圧Vで得るK=\allowbreak{}eVを、一光子へ最大限移すと\(hc/\lambda_{\min}=eV\)である。\(\lambda_{\min}=\frac{6.6\times10^{-34}\times3.0\times10^8}{1.6\times10^{-19}\times35\times10^3}=3.5357\times10^{-11}\mathrm m\)より2。現行h/eの精密値へ置き換えず指定丸め値を使う。電子速度の非相対論近似は不要である。

\Needspace{4\baselineskip}\paragraph{選択肢の検討（査読待ち初稿）}\mbox{}\par

\noindent\textbf{1：不一致} \(2.6\times10^{-11}\mathrm m\)。求めた3.5357e-11mより小さく、本条件の電子エネルギー限界を超える。\par

\noindent\textbf{2：最も近い} \(3.5\times10^{-11}\mathrm m\)。3.5357e-11mへ最も近い提示値。\par

\noindent\textbf{3：不一致} \(3.9\times10^{-11}\mathrm m\)。限界より長い波長の光子は放出され得るが最短値として2より遠い。\par

\noindent\textbf{4：不一致} \(9.1\times10^{-11}\mathrm m\)。連続線にこの波長が存在し得ることと最短端であることは別、算出値から遠い。\par

\noindent\textbf{5：不一致} \(9.9\times10^{-11}\mathrm m\)。この値も最短端ではなく、eVを全て一光子へ与える上限と一致しない。\par

\paragraph{再利用する考え方} 定義の基準日を確定し、物理量の関係から指数と係数を導く。量子現象では運動量・エネルギー保存と観測対象を区別する。

\paragraph{共通理論} \hyperref[theory:modern:xray]{第\ref*{theory:modern:xray}節}, \hyperref[theory:modern:duality]{第\ref*{theory:modern:duality}節}

\par 参考：\cite{os_photoelectric}。

\par 難易度 2（校準初稿）、優先度未評価。

\section{放射線量の物理量・防護量}

\subsection{平成27年度 問6}
\label{q:h27-reg-applied-06}

\noindent\textbf{原典：} sources/H27/17\_applied\_science.pdf、PDF物理3頁。

\textbf{公式正答：} 5（single）。

\paragraph{問題の要旨} 吸収エネルギーと防護上の重み付けを、線量名と単位の対応表へ結び付ける。

\Needspace{4\baselineskip}\paragraph{確認した条件・記述（独自要約）}\mbox{}\par

\noindent\textbf{ア：} 単位質量当たりの吸収エネルギー\par

\noindent\textbf{イ：} アの単位\par

\noindent\textbf{ウ：} 放射線荷重係数をアへ乗じた量（原典の名称）\par

\noindent\textbf{エ：} ウと実効線量の単位\par

\noindent\textbf{方向：} 最も適切な組合せ。\par

\paragraph{問われる核心} エネルギー量と生物影響を考慮する防護量

\paragraph{必要知識} 吸収線量D、Gy=\allowbreak{}J/kg、HT=\allowbreak{}sum wR DTR、E=\allowbreak{}sum wT HT、Sv、Bq=\allowbreak{}s\textasciicircum{}-1、重みは無次元。

\paragraph{考え方・正答を選ぶ理由（査読待ち初稿）} 5である。質量当たり吸収エネルギー\(D=d\epsilon/dm\)の単位はGy=\allowbreak{}J/kg。組織Tへの放射線Rの平均吸収線量に無次元重みwRを掛け\(H_T=\sum_Rw_RD_{T,R}\)とするとSv。さらに\(E=\sum_Tw_TH_T\)が実効線量でSvである。Bqは壊変/秒であってエネルギー/質量ではない。

\Needspace{4\baselineskip}\paragraph{選択肢の検討（査読待ち初稿）}\mbox{}\par

\noindent\textbf{1：単位不適切} 吸収線量／Bq／等価線量／Sv。ア・ウ・エは対応するが吸収線量のイはGyでありBqは壊変率。\par

\noindent\textbf{2：名称と単位不適切} 等価線量／Gy／吸収線量／Bq。アとウが逆である。吸収線量にGy、等価/実効線量にSvを使うのでエBqも誤り。\par

\noindent\textbf{3：単位逆} 吸収線量／Sv／等価線量／Gy。アとウは適切だがGyとSvを逆にしている。次元が同じでも量の意味が違う。\par

\noindent\textbf{4：名称逆} 等価線量／Sv／吸収線量／Gy。Svが等価線量、Gyが吸収線量なのはよいが、原典のアは質量当たりエネルギーなので吸収、ウは放射線重み付け後なので等価。名称を逆にした。\par

\noindent\textbf{5：適切} 吸収線量／Gy／等価線量／Sv。アは吸収線量D、イGy、ウは等価線量HT、エSvである。実効線量EもSvだが組織重み付けをさらに持つ。\par

\paragraph{再利用する考え方} 定義の基準日を確定し、物理量の関係から指数と係数を導く。量子現象では運動量・エネルギー保存と観測対象を区別する。

\paragraph{共通理論} \hyperref[theory:modern:dose]{第\ref*{theory:modern:dose}節}, \hyperref[theory:modern:activity]{第\ref*{theory:modern:activity}節}

\par 参考：\cite{modern_dose,si_brochure9}。

\par 難易度 2（校準初稿）、優先度未評価。
